% Clase 1 --- Relación de pertenencia y diagramas de Venn (§3.1).
% El docente: "en general se dibujan los conjuntos así, con los elementos
% adentro... si A es mi conjunto, tenemos que x pertenece a A, mientras que
% y no pertenece a A". Diagrama de Venn simple: un óvalo (conjunto A) con un
% elemento adentro y uno afuera.
\begin{center}
\begin{tikzpicture}[scale=1]
  % óvalo del conjunto A
  \draw[figset] (0,0) ellipse (2.2 and 1.4);
  \node[figlbl] at (0,1.75) {$A$};

  % elemento x, adentro
  \node[figptoY, minimum size=6pt] (px) at (-0.6,0.25) {};
  \node[figlbl, below=2pt] at (px) {$x$};
  \node[figlbl] at (-0.6,-1.75) {$x \in A$};

  % elemento y, afuera
  \node[figptoZ, minimum size=6pt] (py) at (3.4,0.6) {};
  \node[figlbl, above=2pt] at (py) {$y$};
  \node[figlbl] at (3.4,-1.75) {$y \notin A$};
\end{tikzpicture}\\[0.5em]
{\small\itshape Diagrama de Venn: el conjunto $A$ se dibuja como un óvalo; sus
elementos van adentro ($x$) y los objetos que no pertenecen a $A$ quedan
afuera ($y$).}
\end{center}
